\subsection{基域扩张}

设 \(\mathfrak{g}\) 是 K 上的李代数，\(\mathrm{K}_1\) 是 K 的扩域，\(\mathfrak{g}' = \mathfrak{g}_{(\mathrm{K}_1)}\)。由于 \(\mathcal{C}^k\mathfrak{g}' = (\mathcal{C}^k\mathfrak{g})_{(\mathrm{K}_1)}\)，\(\mathfrak{g}\) 幂零当且仅当 \(\mathfrak{g}'\) 幂零。

设 M 是在 K 上有限维的 g-模，n 是 M 的最大幂零理想，且 \(M' = M_{(K_1)}\)。设 \((\mathrm{M}_i)_{0 \leq i \leq n}\) 是 g-模 M 的 Jordan-Hölder 序列。则对所有 i 及所有 \(x_{\mathrm{M}}(\mathrm{M}_i) \subset \mathrm{M}_{i+1}\)，有 \(x \in n\)，因此

\[
x _ {\mathbf {M} ^ {\prime}} ^ {\prime} \left(\left(\mathrm{M} _ {i}\right) _ {\left(\mathrm{K} _ {1}\right)}\right) \subset \left(\mathrm{M} _ {i + 1}\right) _ {\left(\mathrm{K} _ {1}\right)}
\]

对所有 \(i\) 及所有 \(x' \in \mathfrak{n}_{(\mathbf{K}_1)}\) 成立；因此当 \(x_{\mathbf{M}}'\) 时 \(x' \in \mathfrak{n}_{(\mathbf{K}_1)}\) 是幂零的，从而 \(\mathfrak{n}_{(\mathbf{K}_1)}\) 包含于 \(\mathfrak{n}'\) 的最大幂零理想 \(\mathbf{M}'\)。现在证明：若 \(\mathbf{K}_1\) 在 \(\mathbf{K}\) 上可分，则 \(\mathfrak{n}' = \mathfrak{n}_{(\mathbf{K}_1)}\)。令 \(\mathbf{E}\) 为由 1 和 \(\mathbf{x}_{\mathbf{M}}\)（\(x \in \mathfrak{g}\)）生成的结合 K-代数，E' 为由 1 和 \(\mathbf{x}_{\mathbf{M}}\)（\(x' \in \mathfrak{g}'\)）生成的结合 K-代数，\(\mathbf{R}\) 和 \(\mathbf{R}'\) 分别为 \(\mathbf{E}\) 和 \(\mathbf{E}'\) 的 Jacobson 根。代数 E' 典范地等同于 \(\mathrm{E}_{(\mathbf{K}_1)}\)。于是 \(\mathrm{R}' = \mathrm{R}_{(\mathbf{K}_1)}\)（《代数》，第 VIII 章，§ 7，第 2 号命题 3 的推论 2 (c)）。现在令 \(y' \in \mathfrak{n}'\)，写成 \(y^{\prime} = \sum_{i = 1}^{n}\lambda_{i}y_{i}\)，其中 \(y_{i}\) 属于 \(\mathfrak{g}\)，且 \(\lambda_{i}\in \mathrm{K}_{1}\) 在 K 上线性无关。于是 \(y_{\mathbf{M}'}' = \sum_{i=1}^{n} \lambda_i(y_i)_{\mathbf{M}'}\)，并且 \(y_{\mathbf{M}'}' \in \mathbb{R}' = \mathbb{R}_{(\mathbf{K}_1)}\)。所以 \((y_i)_{\mathbf{M}} \in \mathbb{R}\)，从而对所有 \(y_i \in \mathfrak{n}\) 有 \(i\)。因此 \(y' \in \mathfrak{n}_{(\mathbf{K}_1)}\)，即 \(\mathfrak{n}' \subset \mathfrak{n}_{(\mathbf{K}_1)}\)。

特别地，若 \(\mathbf{K}_1\) 在 \(\mathbf{K}\) 上可分，则 \(\mathfrak{g}_{(\mathbf{K}_1)}\) 的最大幂零理想由 \(\mathfrak{g}\) 的最大幂零理想通过将基域从 \(\mathbf{K}\) 扩张到 \(\mathbf{K}_1\) 而得到。

